\clearpage
\section{Centered sources and graph
estimates}\label{c-centered-sources-and-graph-estimates}

Technical derivation for review. Equation and subsection numbers in this
appendix are local to the source derivation. Historical local labels are
retained, including numbering gaps or nonsequential labels. The
accompanying source notes retain the original expressions for
comparison.

\subsection{0. Result and scope}\label{c-result-and-scope}

For any fixed partition \((n_{1},\ldots ,n_{k})\), the centered
coefficient calculation has an exact root-graph resolution:

\begin{itemize}
\tightlist
\item
  Each two-block charge source occupies two quanta of one root pair and
  has fast energy \(4r_{ab}\).
\item
  Each genuine three-block charge source occupies one quantum on every
  edge of a triangle and has fast energy \(2(r_{ab}+r_{bc}+r_{ca})\).
\item
  On residual block-center-neutral states, the leading inverse-square
  Schur multiplication cancels as an operator on the entire internal
  Clifford module. Internal derivatives of the adapted vacuum are
  included in this statement; freezing the internal coordinates is not
  permission to discard those derivatives.
\item
  The associated finite dressing has pair weight \(C_{N} r_{ab}^{-3}\)
  and triangle weight
  \(C_{N} [r_{ab}^{-2}+r_{bc}^{-2}+r_{ca}^{-2}]/(r_{ab}+r_{bc}+r_{ca})\).
  The latter summed denominator is essential beside a far large block
  and a close opposite pair.
\item
  A graph-local multiplier lemma turns these weights into a small
  multiple of the complete slow supersymmetric form plus a small
  pairwise center-Hardy multiplier, without expanding an inverse around
  an internal Hamiltonian and without applying physical internal Hardy
  to colored vectors.
\end{itemize}

The nonzero internal coefficient continuation is supplied in Appendix
\(D\); analyticity of the oscillator matrices alone would not supply
that continuation.

\subsection{1. Fibers, neutral sector, and the correct fast
inverse}\label{c-fibers-neutral-sector-and-the-correct-fast-inverse}

Write \(b_{ab}=z_{a}-z_{b}\), \(r_{ab}=\vert b_{ab}\vert\), and
\(d_{ab}=n_{a} n_{b}\). Centers have mass metric
\(\sum _{a} n_{a} \vert dz_{a}\vert ^{2}\). Internal matrices
\(A^{(a)}\) are traceless Hermitian, with radii

% DISPLAY c.1
% SOURCE alpha_a² = sum_i ||A_i^(a)||_F².
\begin{gather*}
\begin{aligned}
& \alpha _{a}^{2} = \sum _{i} \Vert A_{i}^{(a)}\Vert _{F}^{2}.
\end{aligned}\notag
\end{gather*}

The hypotheses used in graph estimates are genuinely pairwise:

% DISPLAY c.2
% SOURCE alpha_a <= kappa r_ab  for every b != a,
% SOURCE r_ab >= r_* > 0.
\begin{gather*}
\begin{aligned}
& \alpha _{a} \le  \kappa  r_{ab} \quad  \text{ for every } b \ne  a, \\
& r_{ab} \ge  r_{*} > 0.
\end{aligned}\notag
\end{gather*}

No comparison of one gap with an unrelated gap is made.

At \(A=0\) the fast vacuum is a product
\(U_{0}= \otimes _{a<b} U_{ab}\). Each pair has \(16d_{ab}\) real
transverse bosonic coordinates and \(16d_{ab}\) complex fermionic modes,
with \(8d_{ab}\) occupied modes. The reverse block is the adjoint, not a
second independent pair. The vacuum is invariant under the residual
\(S(\prod  U(n_{a}))\) group and is even in boson parity and fast
fermion parity.

For any one normalized real supercharge let \(q_{0}\) be its centered
fast part. Before imposing the block-center Gauss constraints, its
square contains the center-weighted gauge term:

% DISPLAY c.3
% SOURCE q_0² = H_f + sum_i (Gamma_i)_(alpha alpha)
% SOURCE                      sum_a z_a^i G_(a,center,fast),
\begin{gather*}
\begin{aligned}
& q_{0}^{2} = H_{f} + \sum _{i} (\Gamma _{i})_{\alpha  \alpha} \\
& \sum _{a} z_{a}^{i} G_{a,\mathrm{center},\mathrm{fast}},
\end{aligned}\notag
\end{gather*}

up to the common gauge-generator normalization. The background
\(z_{a} I_{n_{a}}\) annihilates all traceless internal generators in
this term. Consequently

% DISPLAY c.4
% SOURCE q_0² = H_f
\begin{gather*}
\begin{aligned}
& q_{0}^{2} = H_{f}
\end{aligned}\notag
\end{gather*}

on the residual block-center-neutral fast space. This includes fast
states carrying nontrivial internal \(SU(n_{a})\) representations. It
does not assert the identity on arbitrary block-center-charged vectors.
Physical total residual equivariance implies this center neutrality
because traceless internal coordinates and internal adjoint fermions
carry no block-center charge.

Every source below is center neutral: a pair contracts opposite root
charges, while a triangle follows an oriented root cycle. On this space
\(P=U_{0} U_{0}^{*}\) is the kernel of \(H_{f}\) and of \(q_{0}\), and

% DISPLAY c.5
% SOURCE q_0^-1 = q_0 H_f^-1
\begin{gather*}
\begin{aligned}
& q_{0}^{-1} = q_{0} H_{f}^{-1}
\end{aligned}\notag
\end{gather*}

is a genuine inverse on \(1-P\). This is not the inverse of an
interacting internal Hamiltonian.

\subsection{2. The internal vacuum jet cannot be
omitted}\label{c-the-internal-vacuum-jet-cannot-be-omitted}

Let \(U(z,A)\) be the normalized adapted product vacuum from the gapped
quadratic matrices in Appendix \(B\). At one pair, rotate \(n=b/r\) to
direction \(9\) and write

% DISPLAY c.6
% SOURCE T_i = A_i^(a) tensor I - I tensor (A_i^(b))^T,
% SOURCE F = rI + T_9,
% SOURCE D = Gamma_9 r + sum_i Gamma_i T_i.
\begin{gather*}
\begin{aligned}
& T_{i} = A_{i}^{(a)} \otimes  I - I \otimes  (A_{i}^{(b)})^{T}, \\
& F = rI + T_{9}, \\
& D = \Gamma _{9} r + \sum _{i} \Gamma _{i} T_{i}.
\end{aligned}\notag
\end{gather*}

At \(A=0\), the bosonic kinetic correction \(G-I\) has no linear term
and the bosonic frequency has first variation
\(\delta  \Omega _{ij}=\delta _{ij} \delta  T_{9}\). The fermionic
negative projector has first variation

% DISPLAY c.7
% SOURCE delta P_- = -(1/(2r)) sum_(i perpendicular n) Gamma_i tensor delta T_i.
\begin{gather*}
\begin{aligned}
& \delta  P_{-} = -(1/(2r)) \sum _{i \perp  n} \Gamma _{i} \otimes  \delta  T_{i}.
\end{aligned}\notag
\end{gather*}

In a local unitary frame with zero scalar Berry connection at this
point, the vacuum derivative is orthogonal to \(U_{0}\) and obeys the
exact norm identity

% DISPLAY c.8
% SOURCE ||delta U_ab||² = (2/r_ab²) sum_i Tr_(d_ab)(delta Q_i²),       (2.1)
\begin{gather*}
\begin{aligned}
& \Vert \delta  U_{ab}\Vert ^{2} = (2/r_{ab}^{2}) \sum _{i} \operatorname{Tr} _{d_{ab}}(\delta  Q_{i}^{2}),
\end{aligned}\tag{2.1}
\end{gather*}

where \(\delta  Q_{i}=\delta  b_{i} I+\delta  T_{i}\). The longitudinal
contribution is the Gaussian variation; the transverse contributions are
the occupied-determinant variation. Formula (2.1) is a norm in the fast
fiber, leaving the entire slow Clifford module untouched.

For example, summing over all canonical center and endpoint-internal
coordinate directions gives

% DISPLAY c.9
% SOURCE sum_mu ||partial_mu U_ab||²
% SOURCE    = 18 n_a n_b (n_a+n_b) / r_ab².                         (2.2)
\begin{gather*}
\begin{aligned}
& \sum _{\mu} \Vert \partial _{\mu} U_{ab}\Vert ^{2} \\
& = 18 n_{a} n_{b} (n_{a}+n_{b}) / r_{ab}^{2}.
\end{aligned}\tag{2.2}
\end{gather*}

Indeed the center contribution is \(18(n_{a}+n_{b})/r^{2}\), while the
internal contributions are
\(18[n_{b}(n_{a}^{2}-1)+n_{a}(n_{b}^{2}-1)]/r^{2}\). Thus the adapted
vacuum has nonzero internal derivatives even at \(A=0\).

\subsubsection{2.1 Zero first jet of the frozen-charge
defect}\label{c-zero-first-jet-of-the-frozen-charge-defect}

Let \(q_{f}(z,A)\) be the frozen linear fast charge, with the exact
frozen eliminated-Gauss kinetic coefficient. It need not annihilate the
adapted vacuum at nonzero internal commutators. Nevertheless

% DISPLAY c.10
% SOURCE R(z,A)=q_f(z,A)U(z,A),
% SOURCE R(z,0)=0,   partial_A R(z,0)=0.                            (2.3)
\begin{gather*}
\begin{aligned}
& R(z,A)=q_{f}(z,A)U(z,A), \\
& R(z,0)=0, \quad  \partial _{A} R(z,0)=0.
\end{aligned}\tag{2.3}
\end{gather*}

Here is the algebra behind the first derivative. In the frozen
superalgebra the discrepancy between \(q_{f}^{2}\) and the quadratic
Hamiltonian consists of the residual internal gauge term and
background-commutator terms. The latter are quadratic in \(A\). The
derivative of the former annihilates \(U_{0}\), since \(U_{0}\) is
internal-gauge invariant. Differentiating the adapted ground equation
and using \(E_{0}'(0)=0\) therefore gives

% DISPLAY c.11
% SOURCE q_0 [q_f'(0)U_0 + q_0 U'(0)] = 0.
\begin{gather*}
\begin{aligned}
& q_{0} [q_{f}'(0)U_{0} + q_{0} U'(0)] = 0.
\end{aligned}\notag
\end{gather*}

The vector in brackets is center neutral and odd in boson parity: a
linear fast charge maps an even Gaussian to an odd Gaussian-polynomial
vector. The kernel of \(q_{0}\) is the even vacuum line, so the bracket
is zero. This parity step is required; the differentiated square alone
does not rule out a vacuum component.

There is also a direct commuting-axis proof. Vary one canonical internal
coordinate \(A_{i}^{(a)}=t T\), with all other internal matrices zero.
This whole one-parameter background is commuting. A unitary
diagonalizing the single Hermitian \(T\) splits the adapted fast system
into scalar color-pair oscillators at shifted relative vectors. The
kinetic-aware Gaussian is the scalar transverse vacuum written in the
chosen oblique slice, and the filled fermion vacuum is neutral for each
color-pair phase. The frozen charge annihilates this product exactly.
Thus every canonical internal partial derivative of \(R\) at zero
vanishes, and linearity gives the full first jet. This argument is only
along commuting coordinate axes; it does not incorrectly extend
annihilation to noncommuting internal tuples.

Center derivatives of \(R\) vanish as well because \(R(z,0)=0\)
identically in \(z\). Equation (2.3), rather than a false assertion that
the adapted fast charge has an exact kernel for all \(A\), is what the
centered homological calculation uses.

\subsection{3. Exact centered coefficient identity, including internal
derivatives}\label{c-exact-centered-coefficient-identity-including-internal-derivatives}

Under the simultaneous coefficient dilation

% DISPLAY c.12
% SOURCE z=L zeta,  A=L a,  Y=L^(-1/2) xi,
\begin{gather*}
\begin{aligned}
& z=L \zeta , \quad  A=L a, \quad  Y=L^{-1/2} \xi ,
\end{aligned}\notag
\end{gather*}

the exact reduced charge has the formal homogeneous inventory

% DISPLAY c.13
% SOURCE q(L) = L² Q_I^V(a) + L^(1/2) q_0(zeta,a)
% SOURCE        + L^-1 q_1 + L^(-5/2) q_2 + L^-4 q_3 + ... .      (3.1)
\begin{gather*}
\begin{aligned}
& q(L) = L^{2} Q_{I}^{V}(a) + L^{1/2} q_{0}(\zeta ,a) \\
& + L^{-1} q_{1} + L^{-5/2} q_{2} + L^{-4} q_{3} + \ldots  .
\end{aligned}\tag{3.1}
\end{gather*}

This is a coefficient expansion near the zero section, not a claim that
the internal interacting charge is bounded or small. In an intrinsic
multiblock chart the inverse Gauss matrix can create further terms; one
must not copy the finite three-term axial expansion without checking
this. All charge coefficients are first order in slow momenta. Since
\(Q_{I}^{V}(a)\) and its first derivative vanish at \(a=0\), its
anticommutators with \(q_{2}\) and \(q_{3}\) vanish there. Thus the
relevant centered Hamiltonian coefficients are exactly

% DISPLAY c.14
% SOURCE H_1 = {q_0,q_1},
% SOURCE H_2 = q_1² + {q_0,q_2},                                  (3.2)
\begin{gather*}
\begin{aligned}
& H_{1} = \{q_{0},q_{1}\}, \\
& H_{2} = q_{1}^{2} + \{q_{0},q_{2}\},
\end{aligned}\tag{3.2}
\end{gather*}

when evaluated at the centered internal jet. Terms from the quadratic
density of the intrinsic Gauss determinant belong to this inventory. The
fact that its linear-in-Y trace vanishes does not permit omitting its
quadratic triangle terms.

Put

% DISPLAY c.15
% SOURCE D_s = U* q_1 U,
% SOURCE B = (1-P) q_1 U.
\begin{gather*}
\begin{aligned}
& D_{s} = U^{*} q_{1} U, \\
& B = (1-P) q_{1} U.
\end{aligned}\notag
\end{gather*}

The derivative of a coefficient section \(f\) in \(q_{1}(Uf)\) lies in
the vacuum line. Therefore \(B\) is a multiplication map, but it
includes the internal derivatives in (2.1):

% DISPLAY c.16
% SOURCE q_1(Uf) = U D_s f + Bf.
\begin{gather*}
\begin{aligned}
& q_{1}(Uf) = U D_{s} f + Bf.
\end{aligned}\notag
\end{gather*}

Equation (2.3) gives the centered off-diagonal source and diagonal
compression

% DISPLAY c.17
% SOURCE (1-P) H_1 U = q_0 B,
% SOURCE U* H_2 U = D_s² + B*B.                                  (3.3)
\begin{gather*}
\begin{aligned}
& (1-P) H_{1} U = q_{0} B, \\
& U^{*} H_{2} U = D_{s}^{2} + B^{*}B.
\end{aligned}\tag{3.3}
\end{gather*}

These are differential-operator/form identities on the common smooth
coefficient core. There is no projection onto internal bound states or
restriction of the spectator Clifford module.

The exact neutral fast inverse then gives

% DISPLAY c.18
% SOURCE B* q_0 H_f^-1 q_0 B = B*B.                               (3.4)
\begin{gather*}
\begin{aligned}
& B^{*} q_{0} H_{f}^{-1} q_{0} B = B^{*}B.
\end{aligned}\tag{3.4}
\end{gather*}

Consequently the centered effective coefficient is \(D_{s}^{2}\), and
the entire additional multiplication coefficient cancels. It is
important to identify \(D_{s}\) with its first jet before replacing
\(D_{s}^{2}\) by a free slow Laplacian.

\subsubsection{3.1 Compression and its first
jet}\label{c-compression-and-its-first-jet}

The following checks supply that last identification at \(A=0\).

\begin{enumerate}
\def\labelenumi{\arabic{enumi}.}
\item
  The scalar Berry curvature is zero at \(A=0\), including pure internal
  and mixed center/internal components. A local gauge therefore has zero
  connection and zero first jet there.
\item
  The expectation of the quadratic-Y commutator charge is zero to first
  order in \(A\). The centered covariance and its first variation are
  spatial-diagonal, so contraction with \(\Gamma _{ij}\) vanishes.
\item
  The triangle orbital term and the fast-fast spin-Gauss term have an
  odd number of fast fermions and have zero vacuum compression. The same
  remains true after one internal derivative.
\item
  The \(\delta\) part of the mixed slow-fast fermion covariance cancels
  the half-density derivative. The remaining sign-covariance term is, up
  to the fixed charge sign and factors of \(i\), a linear combination of

  % DISPLAY c.19
  % SOURCE Gamma_n Im Tr_color[T_c F^-1 sign(D)],                 (3.5)
  \begin{gather*}
  \begin{aligned}
  & \Gamma _{n} \operatorname{Im}  \operatorname{Tr} _{\mathrm{color}}[T_{c} F^{-1} \operatorname{sign} (D)],
  \end{aligned}\tag{3.5}
  \end{gather*}

  where \(T_{c}\) is the Hermitian bifundamental action of a slow color
  generator. The trace leaves the spin matrix untraced. At the centered
  point it vanishes. Its first variation is a linear combination of

  % DISPLAY c.20
  % SOURCE r^-2 Im Tr(T_c delta T_n) Gamma_n,
  % SOURCE r^-2 Im Tr(T_c delta T_i) Gamma_i.
  \begin{gather*}
  \begin{aligned}
  & r^{-2} \operatorname{Im}  \operatorname{Tr} (T_{c} \delta  T_{n}) \Gamma _{n}, \\
  & r^{-2} \operatorname{Im}  \operatorname{Tr} (T_{c} \delta  T_{i}) \Gamma _{i}.
  \end{aligned}\notag
  \end{gather*}

  Both are zero because the trace of a product of two Hermitian matrices
  is real. This is why the compression must be checked as a jet, not
  only by its value at \(A=0\).
\end{enumerate}

In a realification of the complex pair space, the covariance is \(I/2\)
plus the complex structure times the realification of
\(\operatorname{sign} (D)\), multiplied by \(i/2\). The contracted mixed
Gauss coefficient is the realification of \(F^{-1} T_{c}\). Its
\(\delta\) trace gives the density derivative, and its complex-structure
trace is exactly the imaginary color trace in (3.5). This verifies the
color mechanism without assuming that \(F\) commutes with the internal
generators.

It follows that \(D_{s}\) has the centered free slow Dirac value and
first jet, so the coefficient in (3.3)-(3.4) is the full free
center-plus-internal kinetic operator. At nonzero internal coordinates,
the complete internal potential charge must of course be retained; (3.1)
never authorizes discarding it from an all-vector estimate.

\subsection{4. Root graph grading and exact virtual
denominators}\label{c-root-graph-grading-and-exact-virtual-denominators}

Use a root edge \(e=\{a,b\}\) and a triangle \(t=\{a,b,c\}\). For a root
oscillator let \(N_{e}\) count bosonic excitations plus fermionic
particles and holes relative to its vacuum. At the centered point

% DISPLAY c.21
% SOURCE H_f = 2 sum_e r_e N_e.
\begin{gather*}
\begin{aligned}
& H_{f} = 2 \sum _{e} r_{e} N_{e}.
\end{aligned}\notag
\end{gather*}

The leading complementary charge map has the decomposition

% DISPLAY c.22
% SOURCE B = sum_e B_e + sum_t B_t.                               (4.1)
\begin{gather*}
\begin{aligned}
& B = \sum _{e} B_{e} + \sum _{t} B_{t}.
\end{aligned}\tag{4.1}
\end{gather*}

Every \(B_{e}\) has \(N_{e}=2\) and every other \(N_{f}=0\). Its terms
are the Gaussian and determinant derivatives, the two-fast-variable
commutator with slow fermion output, and the complementary mixed
spin-Gauss contraction. Removing \(P\) removes their only possible
zero-quantum part.

Every \(B_{t}\) has one excitation on each of its three edges. Its
possible charge monomials are exactly of these types, with color
contractions along an oriented cycle:

% DISPLAY c.23
% SOURCE Y_e Y_f theta_g,
% SOURCE r_g^-1 Y_e p_f theta_g,
% SOURCE r_h^-1 theta_e theta_f theta_g.                          (4.2)
\begin{gather*}
\begin{aligned}
& Y_{e} Y_{f} \theta _{g}, \\
& r_{g}^{-1} Y_{e} p_{f} \theta _{g}, \\
& r_{h}^{-1} \theta _{e} \theta _{f} \theta _{g}.
\end{aligned}\tag{4.2}
\end{gather*}

Here \(e,f,g\) are the three distinct triangle edges, while \(h\) is one
of them. A fast Majorana applied once to the filled vacuum creates one
particle or hole. Different edges have independent Gaussian and Fock
factors, so no within-edge vacuum contraction occurs in a triangle
monomial.

Thus

% DISPLAY c.24
% SOURCE H_f B_e = 4r_e B_e,
% SOURCE H_f B_t = 2s_t B_t,       s_t=sum_(e in t) r_e.           (4.3)
\begin{gather*}
\begin{aligned}
& H_{f} B_{e} = 4r_{e} B_{e}, \\
& H_{f} B_{t} = 2s_{t} B_{t}, \quad  s_{t}=\sum _{e \in  t} r_{e}.
\end{aligned}\tag{4.3}
\end{gather*}

Distinct graph summands are orthogonal as fast-space maps. Pair sources
are nonvacuum on only their own edge. Triangle sources are odd in
combined excitation parity on exactly their three edges. The centered
charge preserves each combined excitation parity and annihilates vacuum
factors, so the orthogonality survives application of \(q_{0}\) and
\(H_{f}^{-1}\).

The finite dressing has the explicit graph resolution

% DISPLAY c.25
% SOURCE Z_e = -(4r_e)^-1 q_0 B_e,
% SOURCE Z_t = -(2s_t)^-1 q_0 B_t.                               (4.4)
\begin{gather*}
\begin{aligned}
& Z_{e} = -(4r_{e})^{-1} q_{0} B_{e}, \\
& Z_{t} = -(2s_{t})^{-1} q_{0} B_{t}.
\end{aligned}\tag{4.4}
\end{gather*}

In particular, as positive operators on the full slow Clifford module,

% DISPLAY c.26
% SOURCE Z_e* Z_e = (4r_e)^-1 B_e*B_e,
% SOURCE Z_t* Z_t = (2s_t)^-1 B_t*B_t.                           (4.5)
\begin{gather*}
\begin{aligned}
& Z_{e}^{*} Z_{e} = (4r_{e})^{-1} B_{e}^{*}B_{e}, \\
& Z_{t}^{*} Z_{t} = (2s_{t})^{-1} B_{t}^{*}B_{t}.
\end{aligned}\tag{4.5}
\end{gather*}

Replacing the second denominator by one arbitrarily selected root
frequency is incorrect. Replacing every denominator by the unrelated
global minimum throws away exactly the information needed for large
far-away blocks.

\subsection{5. Explicit coefficient seminorms and shape-independent
graph
estimates}\label{c-explicit-coefficient-seminorms-and-shape-independent-graph-estimates}

All constants below are finite at fixed \(N\). They can be calculated
from the normalized Lie structure constants, \(\gamma\) matrices, and
the displayed source monomials; no compact global gap-ratio chart occurs
in their definition.

One precise conservative convention is the following. Write each pair
source as a sum of elementary maps after factoring \(1/r_{e}\); evaluate
the elementary Hermite/Fock monomials on unit-frequency normalized vacua
and include their operator norms and absolute scalar coefficients in a
sum \(K_{e}\). Include the derivative contribution using (2.1). For a
triangle, group the monomials (4.2) into the three types, factoring
respectively

% DISPLAY c.27
% SOURCE (r_e r_f)^-1/2,
% SOURCE r_g^-1 (r_f/r_e)^1/2,
% SOURCE r_h^-1.
\begin{gather*}
\begin{aligned}
& (r_{e} r_{f})^{-1/2}, \\
& r_{g}^{-1} (r_{f}/r_{e})^{1/2}, \\
& r_{h}^{-1}.
\end{aligned}\notag
\end{gather*}

Let \(K_{t}\) be the sum of their unit-frequency Hermite/Fock norms
times absolute coefficients. \(\Gamma\) matrices have norm one, a real
Majorana has norm \(1/\sqrt{2}\), and degree-two normalized Gaussian
monomials have norms at most \(\sqrt{3}/2\). The remaining coefficients
are normalized color brackets and unit spatial-direction contractions.
This is a finite, explicit coefficient seminorm, not an unspecified
supremum over gap ratios. For example (2.2) and Cauchy-Schwarz give a
derivative-source contribution no larger than

% DISPLAY c.28
% SOURCE K_(e,derivative)² <= 162 N² n_a n_b(n_a+n_b).
\begin{gather*}
\begin{aligned}
& K_{e,\mathrm{derivative}}^{2} \le  162 N^{2} n_{a} n_{b}(n_{a}+n_{b}).
\end{aligned}\notag
\end{gather*}

Enlarge \(K_{t}\), if necessary, by a factor two to account for the
following weight reduction. Triangle geometry gives
\(r_{f}\le r_{e}+r_{g}\), and hence

% DISPLAY c.29
% SOURCE r_g^-1 sqrt(r_f/r_e)
% SOURCE    <= r_g^-1 + (r_e r_g)^-1/2
% SOURCE    <= (3/2)r_g^-1+(1/2)r_e^-1.
\begin{gather*}
\begin{aligned}
& r_{g}^{-1} \sqrt{r_{f}/r_{e}} \\
& \le  r_{g}^{-1} + (r_{e} r_{g})^{-1/2} \\
& \le  (3/2)r_{g}^{-1}+(1/2)r_{e}^{-1}.
\end{aligned}\notag
\end{gather*}

Also \((r_{e} r_{f})^{-1/2} \le  (r_{e}^{-1}+r_{f}^{-1})/2\). Therefore

% DISPLAY c.30
% SOURCE ||B_e|| <= K_e/r_e,
% SOURCE ||B_t|| <= K_t sum_(e in t) r_e^-1.                     (5.1)
\begin{gather*}
\begin{aligned}
& \Vert B_{e}\Vert  \le  K_{e}/r_{e}, \\
& \Vert B_{t}\Vert  \le  K_{t} \sum _{e \in  t} r_{e}^{-1}.
\end{aligned}\tag{5.1}
\end{gather*}

Set \(w_{t}=\sum _{e \in  t} r_{e}^{-2}\). Equations (4.5) and (5.1)
imply

% DISPLAY c.31
% SOURCE ||Z_e||² <= K_e²/(4r_e³),
% SOURCE ||Z_t||² <= 3 K_t² w_t/(2s_t).                          (5.2)
\begin{gather*}
\begin{aligned}
& \Vert Z_{e}\Vert ^{2} \le  K_{e}^{2}/(4r_{e}^{3}), \\
& \Vert Z_{t}\Vert ^{2} \le  3 K_{t}^{2} w_{t}/(2s_{t}).
\end{aligned}\tag{5.2}
\end{gather*}

The centered map by itself has only center derivatives; an internal
derivative requires an extension. For the following internal derivative
bounds, choose smooth gapped pairwise Gaussian/Fock transport of the
centered graph coefficient from \(A=0\), for example transport along
internal radial rays. This gives a concrete graph-factorized extension,
but it is not asserted to equal the actual nonzero-internal homological
source. The center derivative statement is unconditional. For these
derivatives, the same graph coefficients and one differentiated vacuum
factor give, with another computable coefficient seminorm \(K'_{g}\),

% DISPLAY c.32
% SOURCE ||nabla Z_e||² <= K'_e² r_e^-5,
% SOURCE ||nabla Z_t||² <= K'_t² (sum_(e in t) r_e^-1)² w_t/s_t. (5.3)
\begin{gather*}
\begin{aligned}
& \Vert \nabla  Z_{e}\Vert ^{2} \le  {K'_{e}}^{2} r_{e}^{-5}, \\
& \Vert \nabla  Z_{t}\Vert ^{2} \le  {K'_{t}}^{2} (\sum _{e \in  t} r_{e}^{-1})^{2} w_{t}/s_{t}.
\end{aligned}\tag{5.3}
\end{gather*}

In (5.3), \(\nabla\) differentiates the active graph coefficient and its
active oscillator factors. The full map also contains a vacuum factor on
every spectator edge. Its full derivative has the corrected bound

% DISPLAY c.33
% SOURCE ||nabla Z_g||² <= C_N[d_(g,active)+m_g W_2],
% SOURCE W_p=sum_e r_e^-p,                                       (5.4)
\begin{gather*}
\begin{aligned}
& \Vert \nabla  Z_{g}\Vert ^{2} \le  C_{N}[d_{g,\mathrm{active}}+m_{g} W_{2}], \\
& W_{p}=\sum _{e} r_{e}^{-p},
\end{aligned}\tag{5.4}
\end{gather*}

where \(m_{g}\) is the right side of (5.2). The extra term is necessary
when a spectator edge is much shorter. For adapted vacua it includes
internal derivatives of spectator factors as well as center derivatives.
The active triangle right side in (5.3) is at most a fixed numerical
multiple of \({K'_{t}}^{2} \sum _{e \in  t} r_{e}^{-5}\), by Holder and
\(s_{t}\ge \max  r_{e}\). Moreover

% DISPLAY c.34
% SOURCE W_3 W_2 <= (#edges) W_5.
\begin{gather*}
\begin{aligned}
& W_{3} W_{2} \le  (\#\text{edges}) W_{5}.
\end{aligned}\notag
\end{gather*}

Thus the full derivatives still have a summed \(C_{N} W_{5}\) majorant.
Graph-resolved weights must nevertheless be kept when an internal Yukawa
multiplier is to be paid.

\subsubsection{5.1 Perturbed finite source
spaces}\label{c-perturbed-finite-source-spaces}

The Gaussian/Fock statement also has a useful conditional continuation.
On the adapted pairwise-relative tube, unitary/symplectic transport
identifies the finite degree-at-most-three source spaces with their
centered spaces. Frequencies have fixed positive bounds
\(c r_{e} \le  \omega _{e,j} \le  C r_{e}\). A source odd in each
triangle edge therefore has complementary excitation energy at least
\(c s_{t}\), even when modes within one pair split or coincide.

If the exact leading Hamiltonian graph source \(S_{g}(A)\) has the
coefficient bounds

% DISPLAY c.35
% SOURCE ||S_e(A)||² <= C_N/r_e,
% SOURCE ||S_t(A)||² <= C_N s_t w_t,
\begin{gather*}
\begin{aligned}
& \Vert S_{e}(A)\Vert ^{2} \le  C_{N}/r_{e}, \\
& \Vert S_{t}(A)\Vert ^{2} \le  C_{N} s_{t} w_{t},
\end{aligned}\notag
\end{gather*}

then the inverse of the shifted excitation operator

% DISPLAY c.36
% SOURCE H_exc = H_fast - E_(0,total)
\begin{gather*}
\begin{aligned}
& H_{\mathrm{exc}} = H_{\mathrm{fast}} - E_{0,\mathrm{total}}
\end{aligned}\notag
\end{gather*}

on the relevant excited sector satisfies precisely the graph dressing
bounds (5.2), with enlarged constants. The vacuum-energy defect is
handled separately by the general frozen-energy lemma. An unshifted
inverse of the full fast Hamiltonian need not have these denominators,
because unrelated pair vacuum-energy defects would enter it. This is
solely the finite-source fast excitation inverse. No internal momentum
appears in its denominator.

To deduce a small adapted inverse-square remainder from the centered
identity, one additionally needs the exact diagonal \(H_{2}\) inventory,
including its density, ordering, and internal-potential cross terms.
Under coefficientwise perturbation bounds of order \(\kappa\), the
finite resolvent identity gives a defect at most
\(C_{N} \kappa  W_{2}\). Establishing all those hypotheses from the full
reduced Hamiltonian is separate from the graph calculation.

\subsection{6. A graph-local complete slow form
lemma}\label{c-a-graph-local-complete-slow-form-lemma}

This is an all-vector statement, not a finite-Gaussian approximation to
an internal wavefunction.

Let

% DISPLAY c.37
% SOURCE H_s = T_cent + sum_a H_a,
% SOURCE H_a = T_a + V_a + Y_a,
% SOURCE V_a >= 0,     ||Y_a(A)|| <= c_(n_a) alpha_a.
\begin{gather*}
\begin{aligned}
& H_{s} = T_{\mathrm{cent}} + \sum _{a} H_{a}, \\
& H_{a} = T_{a} + V_{a} + Y_{a}, \\
& V_{a} \ge  0, \quad  \Vert Y_{a}(A)\Vert  \le  c_{n_{a}} \alpha _{a}.
\end{aligned}\notag
\end{gather*}

The averaged sixteen-charge identity makes \(H_{a}\) nonnegative also on
colored internal vectors. The ordinary physical lower-rank Hardy
estimate is not being invoked. For a graph-local multiplication map
\(Z_{g}\), suppose

\begin{itemize}
\tightlist
\item
  it is an active-graph coefficient, depending on internal coordinates
  at graph vertices only, tensored with the vacua on spectator edges;
\item
  it acts on internal Clifford factors only at these vertices;
\item
  it is even in total fermion parity, as is \(q_{0}^{-1} B_{g}\);
\item
  \(\Vert Z_{g}\Vert ^{2} \le  m_{g}(z)\) and the full derivative obeys
  the center-only majorant \(d_{g}\) in (5.4).
\end{itemize}

For a block \(a\) outside \(g\), its internal potential charge commutes
exactly with the map. The complete internal charge satisfies

% DISPLAY c.38
% SOURCE Q_(a,alpha) Z_g = Z_g Q_(a,alpha)+C_(a,alpha,g),
\begin{gather*}
\begin{aligned}
& Q_{a,\alpha} Z_{g} = Z_{g} Q_{a,\alpha}+C_{a,\alpha ,g},
\end{aligned}\notag
\end{gather*}

where \(C_{a,\alpha ,g}\) is bounded multiplication coming only from
derivatives of spectator vacua. Its squared norm is bounded by the
corresponding derivative majorant, with the finite Clifford contraction
factor. Averaging the exact charge squares and applying Young gives

% DISPLAY c.39
% SOURCE H_a[Z_g f] <= 2 integral m_g(z) [H_a-density of f]
% SOURCE               + C_N integral d_(a,g)(z)|f|².            (6.1)
\begin{gather*}
\begin{aligned}
& H_{a}[Z_{g} f] \le  2 \int  m_{g}(z) [H_{a}\text{-density of }f] \\
& + C_{N} \int  d_{a,g}(z)\vert f\vert ^{2}.
\end{aligned}\tag{6.1}
\end{gather*}

In particular no unrelated \(\alpha _{a} m_{g}\) loss appears. It would
be incorrect to assert exact commutation while differentiating a moving
spectator vacuum. For a graph vertex use the elementary uncompressed
estimate, with no internal spectral assumptions,

% DISPLAY c.40
% SOURCE H_a[Z_g f]
% SOURCE   <= 2 integral m_g(z) [H_a-density of f]
% SOURCE      + 3c_(n_a) integral alpha_a m_g(z)|f|²
% SOURCE      + 2 integral |partial_Aa Z_g|² |f|².               (6.2)
\begin{gather*}
\begin{aligned}
& H_{a}[Z_{g} f] \\
& \le  2 \int  m_{g}(z) [H_{a}\text{-density of }f] \\
& + 3c_{n_{a}} \int  \alpha _{a} m_{g}(z)\vert f\vert ^{2} \\
& + 2 \int  \vert \partial _{Aa} Z_{g}\vert ^{2} \vert f\vert ^{2}.
\end{aligned}\tag{6.2}
\end{gather*}

The notation for a weighted internal density is harmless here because
its weight depends only on centers; it is the nonnegative internal form
on every fixed-center fiber. The center kinetic estimate has the same
derivative-product bound with no Yukawa term.

For pair and triangle dressings from (5.2),

% DISPLAY c.41
% SOURCE sup m_g <= C_N r_*^-3.
\begin{gather*}
\begin{aligned}
& \sup  m_{g} \le  C_{N} r_{*}^{-3}.
\end{aligned}\notag
\end{gather*}

At a pair endpoint, \(\alpha _{a}\le \kappa  r_{e}\) gives

% DISPLAY c.42
% SOURCE alpha_a m_e <= C_N kappa r_e^-2.
\begin{gather*}
\begin{aligned}
& \alpha _{a} m_{e} \le  C_{N} \kappa  r_{e}^{-2}.
\end{aligned}\notag
\end{gather*}

At any triangle vertex,
\(\alpha _{a}\le \kappa  \min (\text{ incident } r_{e})\le \kappa  s_{t}\),
and hence

% DISPLAY c.43
% SOURCE alpha_a m_t <= C_N kappa w_t.                           (6.3)
\begin{gather*}
\begin{aligned}
& \alpha _{a} m_{t} \le  C_{N} \kappa  w_{t}.
\end{aligned}\tag{6.3}
\end{gather*}

This remains true when the opposite edge is arbitrarily shorter than the
two incident edges. The factor \(s_{t}\) from the virtual denominator is
doing essential work.

Combining (6.1)-(6.3), then using positivity of \(H_{s}\) to sum the
finitely many graph maps, proves

% DISPLAY c.44
% SOURCE H_s[sum_g Z_g f]
% SOURCE   <= C_N r_*^-3 H_s[f]
% SOURCE      + C_N kappa integral W_2 |f|²
% SOURCE      + C_N integral W_5 |f|².                           (6.4)
\begin{gather*}
\begin{aligned}
& H_{s}[\sum _{g} Z_{g} f] \\
& \le  C_{N} r_{*}^{-3} H_{s}[f] \\
& + C_{N} \kappa  \int  W_{2} \vert f\vert ^{2} \\
& + C_{N} \int  W_{5} \vert f\vert ^{2}.
\end{aligned}\tag{6.4}
\end{gather*}

The constants can be chosen from the finite sums of \(K_{g}^{2}\),
\({K'_{g}}^{2}\), vertex counts, and \(c_{n_{a}}\). For instance a
direct Cauchy-Schwarz sum multiplies their sum by the number of graphs
\(\binom{k}{2}+\binom{k}{3}\). This gives an explicit finite-N recipe,
without a global ratio parameter.

The full derivative terms in this argument include spectator factors;
their product weight is controlled by
\(W_{3} W_{2} \le  (\#\text{edges}) W_{5}\). The same argument applies
to a genuinely adapted dressing if its active-graph coefficients and
spectator factors satisfy the hypotheses. In particular, (6.4) retains
arbitrary internal momenta. No absolute-value inverse expansion of
\(H_{f}+H_{I}\) has been used. It also explains why the coarse estimate
\(\Vert Z\Vert ^{2}\le C W_{3}\) should not first be multiplied by
\(\sum  \alpha _{a}\): that would incorrectly charge a distant large
block against an unrelated short root.

\subsubsection{6.1 Direct center-Hardy
payment}\label{c-direct-center-hardy-payment}

For each pair, in the mass center metric,

% DISPLAY c.45
% SOURCE integral r_ab^-2 |f|²
% SOURCE    <= (4/49)(1/n_a+1/n_b)^-1 T_cent[f].                 (6.5)
\begin{gather*}
\begin{aligned}
& \int  r_{ab}^{-2} \vert f\vert ^{2} \\
& \le  (4/49)(1/n_{a}+1/n_{b})^{-1} T_{\mathrm{cent}}[f].
\end{aligned}\tag{6.5}
\end{gather*}

The inequality holds for Hilbert-valued coefficients and unitary
connections. Also \(W_{5} \le  r_{*}^{-3} W_{2}\). Hence the last two
terms in (6.4) are bounded by an arbitrarily small center kinetic
fraction by first choosing \(\kappa\) small at fixed \(N\) and then
\(r_{*}\) large. The first term is a small fraction of the entire
interacting slow form, not a fixed internal kinetic loss.
